\documentclass[border=10pt]{standalone}
\usepackage[utf8]{inputenc}
\usepackage{amsmath}
\usepackage{circuitikz}

\begin{document}
\begin{circuitikz}[european, scale=0.9, font=\small]

% === Source ===
\draw (0, 0.545) to[sV, l=$V_{\mathrm{IN}}$] (0, -4) node[ground]{};
\draw (0, 0.545) to[R, l={RS $100\,\Omega$}] (2, 0.545);

% === Op-amp: yscale=-1 so (+)=upper-left (signal in), (-)=lower-left (feedback) ===
\node[op amp, yscale=-1, anchor=out] (opamp) at (10, 0) {};
\draw (2, 0.545) -| (opamp.+);

% === Inverting input wire: left of op-amp body, then down-right to hub ===
\draw (opamp.-) -- (7.35, -0.545) -- (7.35, -3) -- (11, -3);

% === Output bus (node 3): opamp.out -> C1 tap -> pot CCW -> RL/Out ===
\draw (opamp.out) -- (15, 0);
\node[circ] at (11, 0) {};       % T-junction: C1 top
\node[circ] at (13, 0) {};       % T-junction: pot CCW
\node[circ] at (15, 0) {};       % T-junction: RL, Out
\draw (15, 0) to[R, l={\shortstack[l]{RL\\$100\,\mathrm{k}\Omega$}}] (15, -3.5) node[ground]{};
\draw (15, 0) to[short, -o] (16.5, 0) node[right]{Out};

% === POT1: vertical, CCW(node 3)=top, CW(node 11)=bottom, wiper exits left ===
\draw (13, 0) to[pR, mirror, l={\shortstack[l]{POT1\\$100\,\mathrm{k}\Omega$}}, n=pot1] (13, -6);
% CW indicator: arrow pointing down toward CW end
\draw[-latex, thin] (13, -2.75) -- (13, -3.25);
\node[above, xshift=-0.25cm, font=\bfseries\small] at (pot1.wiper) {Gain};
\node[circ] at (13, -6) {};      % T-junction: pot CW, R1 wire, R2

% === Node 10 hub at (11, -3): wiper, opamp.-, C1 bottom, R1 top ===
\draw (pot1.wiper) -- (11, -3);
\node[circ] at (11, -3) {};      % 4-way junction

% === C1 vertical: parallel with RP1LO (node 3 to node 10) ===
\draw (11, 0) to[C, l={\shortstack[l]{C1\\$2.2\,\mathrm{nF}$}}] (11, -3);

% === R1 vertical: parallel with RP1HI (node 10 to node 11) ===
\draw (11, -3) to[R, l={\shortstack[l]{R1\\$2.2\,\mathrm{k}\Omega$}}] (11, -6);

% Wire: R1 bottom to pot CW end
\draw (11, -6) -- (13, -6);

% === R2 and C2: frequency shaping at CW end ===
\draw (13, -6) to[R, l={\shortstack[l]{R2\\$100\,\Omega$}}] (13, -8.5);
\draw (13, -8.5) to[C, l={\shortstack[l]{C2\\$10\,\mu\mathrm{F}$}}] (13, -10.5) node[ground]{};

\end{circuitikz}
\end{document}
